\section{Related Work}
\label{sec:related}

\para{Involuntary leakage of context.}
Content in a model's context biases its outputs even when it is irrelevant: mentioning that someone likes yellow makes a model more likely to say they drive a school bus~\citep{gonen2025semantic}. \citet{holtzman2026secret} showed that a secret word in the system prompt leaks thematically into stories, with 2AFC discrimination of 64--79\% across frontier models and 83\% for \gemma, and that small models do not leak. Leakage also appears for digit secrets, which can be reconstructed from ordinary outputs~\citep{fairoze2026inadvertent}, and for private information in meeting summaries, where chain-of-thought did not help~\citep{mireshghallah2023confaide}. Negated instructions can raise the probability of the forbidden concept~\citep{mann2025whitebear}, and steering away from a taboo word leaves the concept recoverable~\citep{candussio2026dontsayit}. Theoretical work argues that an agent with only a public transcript cannot keep a secret both hidden and consistent~\citep{baldelli2026hangman}, and implicit secrets drift unless stated~\citep{luo2026beliefs}. Building on \citet{holtzman2026secret}, we ask whether an explicit plan removes the leak and why.

\para{Implicit planning in language models.}
Hidden states encode information about future text. \citet{wu2024planahead} separate pre-caching for future tokens from features that happen to help later. \citet{dong2025emergent} show that prompt activations predict global attributes of a response, such as which animal a story will feature. Paragraph-level plans can be read out at middle-to-late depth~\citep{pochinkov2025parascopes}, and steering a single token can change a planned rhyme~\citep{maar2026plan}. Apparent planning can also reflect plain attention to context~\citep{nainani2025planning}, and attention knockout can isolate lookahead~\citep{men2024lookahead}. These works motivate our hypothesis that an in-context secret shapes implicit plans. Unlike them, we test what happens when the plan is made explicit and supplied from outside.

\para{Explicit plans in story generation.}
Plan-then-write pipelines improve long-story coherence~\citep{yang2022re3,yang2023doc} and suspense~\citep{xie2024suspense}. \citet{sui2026spoiler} show that detailed per-beat plans make late story events less predictable. These systems treat the plan as a way to control quality. We use the plan instead as an intervention on what the writer must decide, and we test whether it stops a held secret from shaping the text.

\para{Reading and removing concepts in activations.}
Secrets fine-tuned into a model can be recovered with the logit lens, sparse autoencoders, and trained activation explainers~\citep{cywinski2025taboo,cywinski2025eliciting,karvonen2025oracles}, although behavioral leakage and decodability can come apart~\citep{bersia2026blindspots}. Fine-tuned model organisms carry constant activation biases that complicate comparisons~\citep{minder2025narrow}, which is one reason we study in-context secrets. We build on difference-in-means directions and directional ablation~\citep{arditi2024refusal}, activation steering~\citep{turner2023actadd,rimsky2023caa}, and concept erasure~\citep{belrose2023leace}. Privacy steering can backfire~\citep{wang2026privacy}, so we compare against controls built in exactly the same way and report text quality.
